\section{Phase and Noise Estimation}
\label{app:frontend}

\subsection{Inter-frame phase estimation}
Stack the received samples as
$[\mathbf Y_{\rm s}]_{qN+r,m}=Y_{qmr}$, so that
$\mathbf Y_{\rm s}\in\mathbb C^{QN\times M}$.
To estimate the phase increments, we approximate the contribution of each
candidate terminal by one dominant path with an unconstrained complex
gain at each antenna.
Absorbing the common scalar phase into the gain, we define the path matrix as
\begin{equation}
 \begin{aligned}
 \mathbf G_0(d,\nu)&=e^{-j2\pi c_1d^2}\mathbf G(d,\nu;0),\\
 [\mathbf h_k(d,\nu)]_{qN+r}
 &=e^{j2\pi\nu qT_{\rm frame}}
 \sum_{n=0}^{N-1}[\mathbf G_0(d,\nu)]_{rn}P_{qkn}.
 \end{aligned}
 \label{eq:phase-atom}
\end{equation}
The delay search is restricted to $d\in[0,N)$ and the Doppler search to $|\nu|\leq\nu_{\max}$. Under $2Nc_1\in\mathbb N_+$, shifting
the delay by an integer multiple of $N$ changes the path operator by
a scalar phase, which the fitted gains absorb. With
$\mathbf h_k=\mathbf h_k(d_k,\nu_k)$ and gain vector
$\mathbf b_k\in\mathbb C^M$, the fitting objective is
\begin{equation}
 \min_{\{d_k,\nu_k,\mathbf b_k\}}
 \left\|\mathbf Y_{\rm s}-\sum_{k=1}^K
 \mathbf h_k\mathbf b_k^{\mathsf T}\right\|_F^2.
 \label{eq:phase-fit}
\end{equation}
When updating terminal $k$, we fix the other contributions and define
$\mathbf Z_k=\mathbf Y_{\rm s}-\sum_{j\ne k}
\mathbf h_j\mathbf b_j^{\mathsf T}$. For a given $(d,\nu)$, the
least-squares gain and the resulting objective are given by
\begin{equation}
 \begin{aligned}
 \mathbf b_k^{\mathsf T}(d,\nu)
 &=\frac{\mathbf h_k(d,\nu)^{\mathsf H}\mathbf Z_k}
 {\|\mathbf h_k(d,\nu)\|^2},\\
 J_k(d,\nu)&=\|\mathbf Z_k\|_F^2-
 \frac{\|\mathbf h_k(d,\nu)^{\mathsf H}\mathbf Z_k\|^2}
 {\|\mathbf h_k(d,\nu)\|^2}.
 \end{aligned}
 \label{eq:phase-variable-projection}
\end{equation}
The gains are eliminated from the separable nonlinear least-squares
problem by variable projection \cite{oleary2013}. Coarse correlations over
delay, within-frame Doppler, and frame phase provide initial candidates.
A phase candidate $\beta$ corresponds to Doppler aliases
\begin{equation}
 \nu_j=\frac{\beta+2\pi j}{2\pi T_{\rm frame}},\qquad
 j\in\mathbb Z,\quad |\nu_j|\leq\nu_{\max}.
 \label{eq:phase-aliases}
\end{equation}
Candidates are ranked using the pilot response vector in
\eqref{eq:phase-atom}, which enforces the same Doppler within and across
frames. The selected candidates undergo bounded local refinement by
L-BFGS-B \cite{zhu1997}, with delay restricted to one differentiable
branch of the wrapped waveform. Each coordinate update is accepted
when it does not increase the residual energy, up to numerical tolerance.
After each sweep, all antenna gains are refitted jointly by least squares.

The receiver uses a fixed number of complete sweeps, retaining the previous
position as an additional refinement starting point after the first sweep.
It sets $\widehat\beta_k=\operatorname{wrap}
(2\pi\widehat\nu_kT_{\rm frame})$ and constructs
$\widehat{\mathbf A}$ by \eqref{eq:effective-pilot}.
This local fit provides phase estimates for all candidate terminals.
The nonlinear denoiser then estimates the full channel tensor in
\eqref{eq:channel}.

\subsection{Noise variance estimation}
Let $L=QN>M\geq2$ and form the spatial sample covariance
\begin{equation}
 \begin{aligned}
 \widehat{\mathbf R}_y&=\frac1L\sum_{q=0}^{Q-1}\sum_{r=0}^{N-1}
 \mathbf y_{qr}\mathbf y_{qr}^{\mathsf H},\\
 \mathbf y_{qr}&=(Y_{q0r},\ldots,Y_{q,M-1,r})^{\mathsf T}.
 \end{aligned}
 \label{eq:spatial-noise-covariance}
\end{equation}
Write its eigenvalues in descending order as
$\lambda_1\geq\cdots\geq\lambda_M$; values below
$\epsilon_{\rm mach}\lambda_1$ are floored at that level before
taking logarithms. For a candidate spatial signal rank $j$, define
\begin{equation}
 \bar\lambda_j=\frac1{M-j}\sum_{i=j+1}^M\lambda_i,\qquad
 \log g_j=\frac1{M-j}\sum_{i=j+1}^M\log\lambda_i.
 \label{eq:noise-tail-statistics}
\end{equation}
The Wax--Kailath minimum description length (MDL) criterion, as reviewed
in \cite{nadakuditi2008}, gives
\begin{equation}
 \begin{aligned}
 \widehat r_{\rm s}
 &=\arg\min_{0\leq j<M}\Big\{
 L(M-j)(\log\bar\lambda_j-\log g_j)\\
 &\hspace{6em}+\frac{j(2M-j)}2\log L\Big\},\\
 \widehat\sigma^2&=\bar\lambda_{\widehat r_{\rm s}}.
 \end{aligned}
 \label{eq:noise-mdl}
\end{equation}
Under a lower-rank spatial signal plus white-noise model, the criterion
selects the spatial signal rank and estimates the noise variance from
the remaining eigenvalues. Both $\widehat\sigma^2$ and
$\widehat{\mathbf A}$ remain fixed throughout the MP iterations.


